\documentclass[aip,preprint]{revtex4-1}

\usepackage[utf8]{inputenc}
\usepackage{amsmath}
\usepackage{amssymb}
\usepackage{braket}
\usepackage{xspace}
\usepackage[capitalize]{cleveref}

\newcommand{\forte}{\textsc{Forte2}\xspace}
\newcommand{\sa}{\mathrm{SA}}
\newcommand{\df}{\mathrm{DF}}
\newcommand{\core}{\mathrm{C}}
\newcommand{\act}{\mathrm{A}}
\newcommand{\virt}{\mathrm{V}}
\newcommand{\Tr}{\operatorname{Tr}}

\begin{document}

\title{Technical Notes: Single-State Nonrelativistic SA-CASSCF Gradients}

\author{Forte2 Developers}

\date{\today}

\maketitle

\section{Scope and assumptions}

This note records the equations implemented in \forte for analytic gradients
of one state evaluated at a state-averaged complete-active-space
self-consistent-field (SA-CASSCF) reference.  The same response equations
apply to SA-GASSCF when the active-space solver imposes the GAS occupation
restrictions and all inter-GAS rotations are optimized.

Only the nonrelativistic implementation is considered.  Its assumptions are
real restricted spatial orbitals, real spin-adapted CI vectors, fixed normalized
state weights, a converged SA-CASSCF or SA-GASSCF reference, and
density-fitted two-electron integrals.  The derivation does not cover complex
orbitals or derivative couplings.  Roots may belong to different spin
multiplicities and therefore to different spin-adapted CI spaces.  Frozen
inactive-core, frozen virtual, and active-frozen orbital response are outside
the current implementation.  Gradients with frozen inter-GAS rotations are
also excluded because defining the displaced frozen orbital frame requires
the response of the orbitals supplied by the parent HF calculation.

\section{Notation}

The real molecular-orbital coefficient matrix satisfies
\begin{equation}
 C^{\mathsf T} S C=I,
\end{equation}
where $S$ is the AO overlap matrix.  The orbitals are partitioned into
inactive core, active, and virtual spaces.  The index conventions are
\begin{center}
\begin{tabular}{llll}
\hline\hline
Space & Symbol & Indices & Dimension \\
\hline
Inactive core & $\mathbb C$ & $m,n,o$ & $N_{\mathrm C}$ \\
Active & $\mathbb A$ & $u,v,w,t$ & $N_{\mathrm A}$ \\
Virtual & $\mathbb V$ & $e,f$ & $N_{\mathrm V}$ \\
Hole, $\mathbb C\cup\mathbb A$ & $\mathbb H$ & $i,j,k,l$ & $N_{\mathrm H}$ \\
General MO & $\mathbb G$ & $p,q,r,s$ & $N_{\mathrm{MO}}$ \\
AO & & $\mu,\nu,\rho,\sigma$ & $N_{\mathrm{AO}}$ \\
Auxiliary AO & & $P,Q,R$ & $N_{\mathrm{aux}}$ \\
\hline\hline
\end{tabular}
\end{center}

For a real CI state $\ket{\Psi_\alpha}=\sum_M c_{M\alpha}\ket{\Phi_M}$,
the active-space spin-free one- and two-particle RDMs are
\begin{align}
 (\gamma^u_v)^\alpha
 &=\sum_\sigma\braket{\Psi_\alpha|
   \hat a^\dagger_{u\sigma}\hat a_{v\sigma}|\Psi_\alpha},
 \\
 (\Gamma^{uv}_{wt})^\alpha
 &=\sum_{\sigma\tau}\braket{\Psi_\alpha|
   \hat a^\dagger_{u\sigma}\hat a^\dagger_{v\tau}
   \hat a_{t\tau}\hat a_{w\sigma}|\Psi_\alpha}.
\end{align}
For normalized weights $w_\alpha$, define
\begin{equation}
 \bar\gamma=\sum_\alpha w_\alpha\gamma^\alpha,
 \qquad
 \bar\Gamma=\sum_\alpha w_\alpha\Gamma^\alpha.
 \label{eq:nr-average-rdms}
\end{equation}
The orbital optimizer permutes the two-particle RDM to the working array
\begin{equation}
 \bar D_{tu,vw}
 =\frac{1}{2}\left(\bar\Gamma^{tv}_{uw}+\bar\Gamma^{uv}_{tw}\right).
 \label{eq:nr-working-rdm}
\end{equation}

The non-antisymmetrized two-electron integrals are in physicists' notation,
\begin{equation}
 \langle pq|rs\rangle
 =\iint\phi_p(1)\phi_q(2)r_{12}^{-1}\phi_r(1)\phi_s(2)
 \,d1\,d2.
 \label{eq:nr-physicists-eri}
\end{equation}
Let the raw three-center integrals and Coulomb metric be
\begin{equation}
 J^P_{\mu\nu}=(P|\mu\nu),
 \qquad M_{PQ}=(P|Q),
 \qquad N=M^{-1}.
 \label{eq:nr-raw-df}
\end{equation}
If small metric eigenvalues are removed, $N$ denotes the pseudoinverse on the
retained auxiliary subspace.  The Fock and response code uses a factor
$\mathsf X$ satisfying
\begin{equation}
 \mathsf X^{\mathsf T}\mathsf X=N,
 \qquad
 B^P_{\mu\nu}=\sum_Q\mathsf X_{PQ}J^Q_{\mu\nu},
 \label{eq:nr-orthonormal-df}
\end{equation}
so that
\begin{equation}
 \sum_P B^P_{\mu\nu}B^P_{\rho\sigma}
 =\sum_{PQ}J^P_{\mu\nu}N_{PQ}J^Q_{\rho\sigma}.
 \label{eq:nr-df-identity}
\end{equation}
With metric truncation, $\mathsf X$ is a symmetric inverse square root on the
retained subspace.  Without truncation, it may instead be an inverse Cholesky
factor.  Response contractions at fixed nuclei use $B$; nuclear derivatives
are evaluated from the raw $J$ and $M$ tensors in \cref{sec:nr-gradient}.

\section{Stationary SA-CASSCF reference}

The CI vectors obey
\begin{equation}
 H_\alpha\mathbf c_\alpha=E_\alpha\mathbf c_\alpha,
 \qquad
 \mathbf c_\alpha^{\mathsf T}\mathbf c_\beta=\delta_{\alpha\beta}
\end{equation}
within each state solver.  The inactive-core and active generalized Fock
contributions are
\begin{align}
 (F_{\core})^q_p
 &=h^q_p+\sum_{m\in\mathbb C}
 \left(2\langle pm|qm\rangle-\langle pm|mq\rangle\right),
 \label{eq:nr-fock-core}
 \\
 (\bar F_{\act})^q_p
 &=\sum_{uv\in\mathbb A}\bar\gamma^u_v
 \left(\langle pu|qv\rangle-\frac12\langle pu|vq\rangle\right),
 \label{eq:nr-fock-active}
 \\
 \bar F^q_p&=(F_{\core})^q_p+(\bar F_{\act})^q_p.
\end{align}
The inactive-core energy is
\begin{equation}
 E_{\core}=\sum_{m\in\mathbb C}
 \left[h^m_m+(F_{\core})^m_m\right].
\end{equation}
The orbital-dependent state-average energy evaluated by the optimizer is
\begin{equation}
 \bar E_{\mathrm{orb}}
 =E_{\core}+V_{\mathrm{NN}}
 +\sum_{uv\in\mathbb A}(F_{\core})^v_u\bar\gamma^u_v
 +\frac12\sum_{tuvw\in\mathbb A}
 \langle tv|uw\rangle\bar D_{tu,vw}.
\end{equation}

\subsection{Orbital Lagrangian and gradient}

The implementation constructs the real orbital Lagrangian matrix $\bar A$
one column space at a time:
\begin{align}
 \bar A^m_p&=2\bar F^m_p,
 &&m\in\mathbb C,
 \label{eq:nr-orbital-lagrangian-core}
 \\
 \bar A^u_p
 &=\sum_{v\in\mathbb A}(F_{\core})^v_p\bar\gamma^v_u
 +\sum_{tvw\in\mathbb A}\langle pv|tw\rangle\bar D_{tu,vw},
 &&u\in\mathbb A,
 \label{eq:nr-orbital-lagrangian-active}
 \\
 \bar A^e_p&=0,
 &&e\in\mathbb V.
\end{align}
Let
\begin{equation}
 \mathbb K=\big((p_I,q_I)\big)_{I=0}^{n_{\mathrm{rot}}-1}
 \label{eq:nr-pair-order}
\end{equation}
be the ordered list selected by the Boolean nonredundant-rotation mask.  The
mask contains exactly one orientation of every retained unordered pair and
is traversed using NumPy C-order Boolean indexing.  The same ordering is used
for every orbital response input and output.  For each pair,
\begin{equation}
 (\bar g)_I
 =2\left(\bar A^{q_I}_{p_I}-\bar A^{p_I}_{q_I}\right).
 \label{eq:nr-orbital-gradient}
\end{equation}
The factor of two is the real restricted spin-free convention used by
\forte.  At convergence, $\bar{\mathbf g}=0$.

For target state $\alpha$, $A^\alpha$ is obtained from
\cref{eq:nr-orbital-lagrangian-core,eq:nr-orbital-lagrangian-active} using
the state-specific RDMs.  Its positive orbital right-hand side is
\begin{equation}
 \boxed{
 (b^\mathrm{o}_\alpha)_I
 =2\left[(A^\alpha)^{q_I}_{p_I}-(A^\alpha)^{p_I}_{q_I}\right]
 }.
 \label{eq:nr-orbital-b}
\end{equation}
No state weight multiplies this target-state quantity, and
\begin{equation}
 \sum_\alpha w_\alpha\mathbf b^\mathrm{o}_\alpha
 =\bar{\mathbf g}=0.
\end{equation}

\subsection{Single-state Lagrangian}

For a target root $\alpha$, the real Lagrangian may be organized as
\begin{align}
 \mathcal L_\alpha
 &=E_\alpha
 -\Tr\!\left[\omega_\alpha(C^{\mathsf T}SC-I)\right]
 +2\sum_\beta\mathbf x_\beta^{\mathsf T}
  (H_\beta-E_\beta I)\mathbf c_\beta
 \nonumber\\
 &\quad+\mathbf z_\alpha^{\mathsf T}\bar{\mathbf g}
 +\text{CI orthonormality constraints}.
 \label{eq:nr-lagrangian}
\end{align}
The orbital multiplier $\mathbf z_\alpha$ and concatenated CI multiplier
$\mathbf x_\alpha$ make this Lagrangian stationary with respect to the
state-averaged orbital and CI parameters.  The CI normalization and
root-rotation constraints are imposed below by projection rather than by
explicitly solving for their multipliers.

\section{Coupled orbital and CI response}

For the ordered pair vector $\mathbf z$, define the antisymmetric embedding
\begin{equation}
 Z_{pq}(\mathbf z)=\sum_Iz_I
 \left(\delta_{p p_I}\delta_{q q_I}
      -\delta_{p q_I}\delta_{q p_I}\right).
 \label{eq:nr-vector-embedding}
\end{equation}
The corresponding orbital path and tangent are
\begin{equation}
 C(\epsilon)=C\exp(\epsilon Z),
 \qquad
 \dot C=\left.\frac{dC(\epsilon)}{d\epsilon}\right|_0=CZ.
 \label{eq:nr-orbital-direction}
\end{equation}

\subsection{Orbital--orbital block}

Define the forward orbital--orbital derivative at fixed SA RDMs and fixed AO
integrals by
\begin{equation}
 \boxed{
 \mathcal A^{\mathrm{oo}}_{IJ}=D_J\bar g_I,
 \qquad
 [\mathcal A^{\mathrm{oo}}\mathbf z]_I
 =\left.\frac{d}{d\epsilon}\bar g_I(Ce^{\epsilon Z})\right|_0
 =2\left[(\dot{\bar A})^{q_I}_{p_I}
         -(\dot{\bar A})^{p_I}_{q_I}\right]
 }.
 \label{eq:nr-oo-action}
\end{equation}
The active RDMs do not change in this partial derivative.  With
$P_{\core}=C_{\core}C_{\core}^{\mathsf T}$ and
$\bar P_{\act}=C_{\act}\bar\gamma C_{\act}^{\mathsf T}$,
\begin{align}
 \dot P_{\core}
 &=\dot C_{\core}C_{\core}^{\mathsf T}
  +C_{\core}\dot C_{\core}^{\mathsf T},
 \\
 \dot{\bar P}_{\act}
 &=\dot C_{\act}\bar\gamma C_{\act}^{\mathsf T}
  +C_{\act}\bar\gamma\dot C_{\act}^{\mathsf T}.
 \label{eq:nr-density-response}
\end{align}
Introduce the linear DF Coulomb and exchange maps
\begin{align}
 J[D]_{\mu\nu}
 &=\sum_{P\rho\sigma}
 B^P_{\mu\nu}B^P_{\rho\sigma}D_{\sigma\rho},
 \\
 K[D]_{\mu\nu}
 &=\sum_{P\rho\sigma}
 B^P_{\mu\sigma}D_{\sigma\rho}B^P_{\nu\rho}.
\end{align}
The AO Fock responses are
\begin{align}
 \dot F_{\core}^{\mathrm{AO}}
 &=2J[\dot P_{\core}]-K[\dot P_{\core}],
 \\
 \dot{\bar F}_{\act}^{\mathrm{AO}}
 &=J[\dot{\bar P}_{\act}]
  -\frac12K[\dot{\bar P}_{\act}].
 \label{eq:nr-fock-response}
\end{align}
For $X\in\{\core,\act\}$, transformation to the moving MO frame gives
\begin{equation}
 \dot F_X=Z^{\mathsf T}F_X+F_XZ
 +C^{\mathsf T}\dot F_X^{\mathrm{AO}}C.
 \label{eq:nr-mo-fock-response}
\end{equation}

The implementation evaluates exchange responses from low-rank factors.  For
the generic symmetric response density
\begin{equation}
 D=C_LG_sC_R^{\mathsf T}+C_RG_sC_L^{\mathsf T},
 \qquad G_s=\frac12(G+G^{\mathsf T}),
 \label{eq:nr-low-rank-density}
\end{equation}
use the signed eigendecomposition
\begin{align}
 G_s&=U\Xi U^{\mathsf T},
 &L&=C_LU|\Xi|^{1/2},
 &R&=C_RU\operatorname{sgn}(\Xi)|\Xi|^{1/2}.
 \label{eq:nr-signed-factorization}
\end{align}
Then $D=LR^{\mathsf T}+RL^{\mathsf T}$.  With half transforms
\begin{equation}
 Y^{P,L}_{\mu a}=\sum_\nu B^P_{\mu\nu}L_{\nu a},
 \qquad
 Y^{P,R}_{\mu a}=\sum_\nu B^P_{\mu\nu}R_{\nu a},
\end{equation}
the exact exchange response is
\begin{equation}
 K[D]_{\mu\nu}=\sum_{Pa}
 \left(Y^{P,L}_{\mu a}Y^{P,R}_{\nu a}
      +Y^{P,R}_{\mu a}Y^{P,L}_{\nu a}\right).
 \label{eq:nr-low-rank-exchange}
\end{equation}
The core and active tuples $(C_L,C_R,G;c_J,c_K)$ are, respectively,
\begin{equation}
 (\dot C_{\core},C_{\core},I;2,1),
 \qquad
 (\dot C_{\act},C_{\act},\bar\gamma;1,1/2),
\end{equation}
where the returned Fock response is $c_JJ[D]-c_KK[D]$.  The directional
exchange workspace is
$O(N_{\mathrm{aux}}N_{\mathrm{AO}}r)$ for factor rank $r$.

Only the general--active transformed DF block is persistent:
\begin{equation}
 B^P_{pu}=\sum_{\mu\nu}C_{\mu p}B^P_{\mu\nu}C_{\nu u},
 \qquad p\in\mathbb G,\quad u\in\mathbb A.
 \label{eq:nr-general-active-df}
\end{equation}
Its response is evaluated without a general--general DF tensor,
\begin{equation}
 \dot B^P_{pu}
 =\sum_rZ_{rp}B^P_{ru}
 +\sum_{\mu\nu}C_{\mu p}B^P_{\mu\nu}\dot C_{\nu u}.
 \label{eq:nr-general-active-df-response}
\end{equation}
The active two-particle contraction is
\begin{align}
 T^P_{tu}&=\sum_{vw\in\mathbb A}B^P_{vw}\bar D_{tu,vw},
 &\dot T^P_{tu}&=\sum_{vw\in\mathbb A}\dot B^P_{vw}\bar D_{tu,vw},
 \label{eq:nr-df-working-rdm}
 \\
 \left.\dot{\bar A}^{u}_{p}\right|_{2e}
 &=\sum_{Pt}\left(\dot B^P_{pt}T^P_{tu}
                  +B^P_{pt}\dot T^P_{tu}\right).
 \label{eq:nr-df-lagrangian-response}
\end{align}
Together with the one-electron/Fock terms, the complete Lagrangian response is
\begin{align}
 \dot{\bar A}^{m}_{p}
 &=2(\dot F_{\core}+\dot{\bar F}_{\act})^{m}_{p},
 &&m\in\mathbb C,
 \\
 \dot{\bar A}^{u}_{p}
 &=\sum_{v\in\mathbb A}(\dot F_{\core})^{v}_{p}\bar\gamma^{v}_{u}
  +\left.\dot{\bar A}^{u}_{p}\right|_{2e},
 &&u\in\mathbb A,
 \label{eq:nr-lagrangian-response}
 \\
 \dot{\bar A}^{e}_{p}&=0,
 &&e\in\mathbb V.
\end{align}
The matrix elements are defined by the matrix-free action,
\begin{equation}
 (\mathcal A^{\mathrm{oo}})_{IJ}
 =[\mathcal A^{\mathrm{oo}}\mathbf e_J]_I.
 \label{eq:nr-oo-dense}
\end{equation}
No four-index integral with general orbital indices is formed, and the
persistent transformed response workspace is
$O(N_{\mathrm{aux}}N_{\mathrm{MO}}N_{\mathrm A})$.

\subsection{Orbital--CI block}

Let the concatenated CI multiplier contain one block for every absolute root,
\begin{equation}
 \mathbf x=(\mathbf x_0^{\mathsf T},\mathbf x_1^{\mathsf T},\ldots,
 \mathbf x_{d-1}^{\mathsf T})^{\mathsf T}.
 \label{eq:nr-ci-layout}
\end{equation}
Blocks occur in absolute-root order and use the CSF ordering of the associated
state solver.  They need not all have the same dimension.  In particular, a
single state average may combine different spin multiplicities, whose
spin-adapted CSF spaces and sigma builders have different dimensions; every
root block is dispatched to its own state solver.

For arbitrary real vectors $\mathbf a$ and $\mathbf b$ in the same CSF space,
define
\begin{align}
 S^{\mathbf a,\mathbf b}&=\mathbf a^{\mathsf T}\mathbf b,
 \\
 (\gamma^u_v)^{\mathbf a,\mathbf b}
 &=\sum_{MN}a_M\braket{\Phi_M|\hat E^u_v|\Phi_N}b_N,
 \\
 (\Gamma^{uv}_{wt})^{\mathbf a,\mathbf b}
 &=\sum_{MN}a_M\braket{\Phi_M|\hat E^{uv}_{wt}|\Phi_N}b_N,
\end{align}
where
\begin{equation}
 \hat E^u_v=\sum_\sigma\hat a^\dagger_{u\sigma}\hat a_{v\sigma},
 \qquad
 \hat E^{uv}_{wt}=\sum_{\sigma\tau}
 \hat a^\dagger_{u\sigma}\hat a^\dagger_{v\tau}
 \hat a_{t\tau}\hat a_{w\sigma}.
\end{equation}
The real bra-plus-ket CI response quantities are
\begin{align}
 s[\mathbf x]
 &=\sum_\beta\left(
 S^{\mathbf x_\beta,\mathbf c_\beta}
 +S^{\mathbf c_\beta,\mathbf x_\beta}\right),
 \label{eq:nr-ci-overlap-response}
 \\
 \gamma^u_v[\mathbf x]
 &=\sum_\beta\left[
 (\gamma^u_v)^{\mathbf x_\beta,\mathbf c_\beta}
 +(\gamma^u_v)^{\mathbf c_\beta,\mathbf x_\beta}\right],
 \\
 \Gamma^{uv}_{wt}[\mathbf x]
 &=\sum_\beta\left[
 (\Gamma^{uv}_{wt})^{\mathbf x_\beta,\mathbf c_\beta}
 +(\Gamma^{uv}_{wt})^{\mathbf c_\beta,\mathbf x_\beta}\right].
 \label{eq:nr-ci-rdm-response}
\end{align}
The working response $D_{tu,vw}[\mathbf x]$ follows from
\cref{eq:nr-working-rdm} with $\bar\Gamma$ replaced by
$\Gamma[\mathbf x]$.  Similarly, $F_{\act}[\mathbf x]$ follows from
\cref{eq:nr-fock-active} with $\bar\gamma$ replaced by
$\gamma[\mathbf x]$.

The CI contribution to the orbital Lagrangian is
\begin{align}
 A^{\mathrm{oc},m}_{p}[\mathbf x]
 &=2\left[s[\mathbf x](F_{\core})^m_p
         +(F_{\act}[\mathbf x])^m_p\right],
 &&m\in\mathbb C,
 \label{eq:nr-oc-lagrangian-core}
 \\
 A^{\mathrm{oc},u}_{p}[\mathbf x]
 &=\sum_{v\in\mathbb A}(F_{\core})^v_p\gamma^v_u[\mathbf x]
 +\sum_{tvw\in\mathbb A}\langle pv|tw\rangle
   D_{tu,vw}[\mathbf x],
 &&u\in\mathbb A,
 \label{eq:nr-oc-lagrangian-active}
 \\
 A^{\mathrm{oc},e}_{p}[\mathbf x]&=0,
 &&e\in\mathbb V.
\end{align}
In the DF implementation, $F_{\act}[\mathbf x]$ uses the low-rank build in
\cref{eq:nr-low-rank-density,eq:nr-low-rank-exchange} with
\begin{equation}
 (C_L,C_R,G;c_J,c_K)
 =(C_{\act},C_{\act},\gamma[\mathbf x]/2;1,1/2).
\end{equation}
The factor $1/2$ in $G$ compensates the two identical terms in
\cref{eq:nr-low-rank-density}.  The two-particle term is contracted as
\begin{equation}
 T^P_{tu}[\mathbf x]
 =\sum_{vw\in\mathbb A}B^P_{vw}D_{tu,vw}[\mathbf x],
 \qquad
 \left.A^{\mathrm{oc},u}_{p}[\mathbf x]\right|_{2e}
 =\sum_{Pt}B^P_{pt}T^P_{tu}[\mathbf x].
 \label{eq:nr-df-oc}
\end{equation}
Thus neither a CI-to-RDM Jacobian nor a general four-index tensor is formed.

The rectangular block action and its dense definition are
\begin{align}
 \boxed{
 [\mathcal A^{\mathrm{oc}}\mathbf x]_I
 =2\left(A^{\mathrm{oc},q_I}_{p_I}[\mathbf x]
         -A^{\mathrm{oc},p_I}_{q_I}[\mathbf x]\right)
 },
 \label{eq:nr-oc-action}
 \\
 (\mathcal A^{\mathrm{oc}})_{IJ}
 &=[\mathcal A^{\mathrm{oc}}\mathbf e_J]_I.
\end{align}
Here $J$ traverses the root-major coefficient layout.  State weights do not
appear in this block because it differentiates the unweighted CI residual in
\cref{eq:nr-lagrangian}.  After projection, $s[\mathbf x]=0$.

\subsection{CI--orbital block}

The transpose mixed block is expressed through the orbital derivative of the
active-space Hamiltonian.  The inactive-core scalar response is
\begin{align}
 \dot E_{\core}[\mathbf z]
 &=\Tr\!\left[\dot C_{\core}^{\mathsf T}
 (h+F_{\core}^{\mathrm{AO}})C_{\core}\right]
 +\Tr\!\left[C_{\core}^{\mathsf T}
 (h+F_{\core}^{\mathrm{AO}})\dot C_{\core}\right]
 \nonumber\\
 &\quad+\Tr\!\left[C_{\core}^{\mathsf T}
 \dot F_{\core}^{\mathrm{AO}}C_{\core}\right].
 \label{eq:nr-core-energy-response}
\end{align}
The nuclear-repulsion energy is fixed in this partial derivative.  Define
\begin{equation}
 \hat H[\mathbf z]
 =\dot E_{\core}[\mathbf z]
 +\sum_{uv\in\mathbb A}(\dot F_{\core})^u_v[\mathbf z]\hat E^u_v
 +\frac12\sum_{uvtw\in\mathbb A}
 \dot{\langle uv|tw\rangle}[\mathbf z]\hat E^{uv}_{tw}.
 \label{eq:nr-response-hamiltonian}
\end{equation}
Only active four-index response integrals are formed:
\begin{align}
 \dot B^P_{uv}
 &=\sum_{p\in\mathbb G}
 \left(Z_{pu}B^P_{pv}+B^P_{up}Z_{pv}\right),
 \\
 \dot{\langle uv|tw\rangle}
 &=\sum_P\left(\dot B^P_{ut}B^P_{vw}
               +B^P_{ut}\dot B^P_{vw}\right).
 \label{eq:nr-active-df-h-response}
\end{align}
The existing CI sigma operation applies $\hat H[\mathbf z]$ to each reference
vector.  For root $\beta$,
\begin{equation}
 \boxed{
 [\mathcal A^{\mathrm{co}}\mathbf z]_\beta
 =2w_\beta\hat H[\mathbf z]\mathbf c_\beta
 }.
 \label{eq:nr-co-action}
\end{equation}
This is the raw CSF-space action, before removing normalization and
root-rotation components.  If $W_{\mathrm C}$ repeats $w_\beta$ for every
coefficient in root block $\beta$, equality of mixed derivatives gives
\begin{equation}
 \mathcal A^{\mathrm{co}}
 =W_{\mathrm C}(\mathcal A^{\mathrm{oc}})^{\mathsf T}.
 \label{eq:nr-mixed-transpose}
\end{equation}

\subsection{CI--CI block and CI projection}

Holding the orbitals and CI energy multipliers fixed, the real CI residual and
its directional derivative are
\begin{align}
 \mathbf r^\mathrm{c}_\beta
 &=2(H_\beta-E_\beta I_\beta)\mathbf c_\beta,
 \\
 \boxed{
 [\mathcal A^{\mathrm{cc}}\mathbf x]_\beta
 =2(H_\beta-E_\beta I_\beta)\mathbf x_\beta
 }.
 \label{eq:nr-cc-action}
\end{align}
The factor of two is the derivative with respect to real CI coordinates.
The dense block is the direct sum
\begin{equation}
 \mathcal A^{\mathrm{cc}}
 =\bigoplus_\beta2(H_\beta-E_\beta I_\beta).
\end{equation}
It is unweighted and singular along every optimized reference vector.

Let $s(\beta)$ identify the state solver containing absolute root $\beta$ and
let $\mathcal R_s$ be its optimized-root set.  The projector in that solver is
\begin{equation}
 Q_s=I_s-\sum_{\gamma\in\mathcal R_s}
 \mathbf c_\gamma\mathbf c_\gamma^{\mathsf T}.
 \label{eq:nr-ci-projector}
\end{equation}
The block-diagonal projector $\mathcal Q_{\mathrm C}$ applies $Q_{s(\beta)}$
to every absolute-root block, and
$\mathcal P_{\mathrm C}=I-\mathcal Q_{\mathrm C}$.

For target root $\alpha$, the raw and projected positive CI right-hand sides
are
\begin{align}
 (\widetilde{\mathbf b}^{\mathrm c}_\alpha)_\beta
 &=2\delta_{\alpha\beta}H_\beta\mathbf c_\beta
 =2\delta_{\alpha\beta}E_\alpha\mathbf c_\alpha,
 \\
 \boxed{
 (\mathbf b^{\mathrm c}_\alpha)_\beta
 =Q_{s(\beta)}(\widetilde{\mathbf b}^{\mathrm c}_\alpha)_\beta=0
 }.
 \label{eq:nr-ci-b}
\end{align}
The code evaluates the raw vector with a sigma operation and explicitly
projects it.  Although the projected CI right-hand side is zero, the coupled
orbital response generally produces a nonzero CI multiplier.

\subsection{Projected coupled solution}

The implemented response equations are
\begin{equation}
 \boxed{
 \begin{pmatrix}
 \mathcal A^{\mathrm{oo}}
 &\mathcal A^{\mathrm{oc}}\mathcal Q_{\mathrm C}\\
 \mathcal Q_{\mathrm C}\mathcal A^{\mathrm{co}}
 &\mathcal Q_{\mathrm C}\mathcal A^{\mathrm{cc}}
  \mathcal Q_{\mathrm C}+\mathcal P_{\mathrm C}
 \end{pmatrix}
 \begin{pmatrix}\mathbf z_\alpha\\\mathbf x_\alpha\end{pmatrix}
 =-
 \begin{pmatrix}\mathbf b^\mathrm{o}_\alpha\\
                 \mathbf b^\mathrm{c}_\alpha\end{pmatrix}
 }.
 \label{eq:nr-response-system}
\end{equation}
The $\mathcal P_{\mathrm C}$ term assigns unit eigenvalues to excluded gauge
directions and enforces
$\mathcal P_{\mathrm C}\mathbf x_\alpha=0$.  Because
\cref{eq:nr-mixed-transpose} contains the coefficient-space weights, the
matrix is generally nonsymmetric in this convention.  \forte applies all four
blocks without forming them and solves \cref{eq:nr-response-system} with
GMRES.  A block-Jacobi preconditioner uses the diagonal orbital Hessian and
diagonal shifted CI Hamiltonians; its storage is linear in the response
dimension.

The pair set $\mathbb K$ is exactly the set optimized in the reference
SA-CASSCF/GASSCF calculation and includes all inter-GAS rotations.

\subsection{Orbital orthogonality multiplier}

The orbital--orbital contribution to the energy-weighted density is not just
$\dot{\bar A}[\mathbf z_\alpha]$, because that derivative is evaluated in a
moving MO frame.  For a general variation $\dot C=CX$,
\begin{equation}
 D_X\left(\mathbf z_\alpha^{\mathsf T}\bar{\mathbf g}\right)
 =2\sum_{pq}X_{pq}
 \left[\dot{\bar A}[\mathbf z_\alpha]
 +Z_\alpha\bar A-\bar A Z_\alpha\right]_{pq}.
 \label{eq:nr-full-coefficient-derivative}
\end{equation}
The commutator is required even though its antisymmetric contribution drops
out of the nonredundant response equation at convergence.  Define
\begin{equation}
 \boxed{
 \Omega_\alpha=A^\alpha
 +A^{\mathrm{oc}}[\mathbf x_\alpha]
 +\dot{\bar A}[\mathbf z_\alpha]
 +Z_\alpha\bar A-\bar A Z_\alpha
 }.
 \label{eq:nr-relaxed-lagrangian}
\end{equation}
The solved response implies
\begin{equation}
 2\left[(\Omega_\alpha)^{q_I}_{p_I}
       -(\Omega_\alpha)^{p_I}_{q_I}\right]=0
 \qquad\text{for every }(p_I,q_I)\in\mathbb K.
\end{equation}
Hermiticity of the orbital orthogonality constraint selects
\begin{equation}
 \boxed{
 \omega_\alpha=\frac12
 (\Omega_\alpha+\Omega_\alpha^{\mathsf T})
 },
 \qquad
 W^{S,\alpha}=C\omega_\alpha C^{\mathsf T}.
 \label{eq:nr-omega}
\end{equation}
Its gradient contribution is
$-\sum_{\mu\nu}W^{S,\alpha}_{\mu\nu}S^x_{\mu\nu}$.

\section{Relaxed densities and DF gradient assembly}
\label{sec:nr-gradient}

Let $x$ denote a Cartesian nuclear displacement and let a superscript $x$
denote a skeleton AO-integral derivative.  The target-root gradient requires
the target densities, the CI-multiplier transition densities, and the orbital
response of the state-average densities.

For fixed full-MO RDMs, define the orbital-frame response generated by
$Z_\alpha$ as
\begin{align}
 \mathcal R_{Z_\alpha}(\bar\gamma)_{pq}
 &=\sum_t\left[(Z_\alpha)_{pt}\bar\gamma_{tq}
              +(Z_\alpha)_{qt}\bar\gamma_{pt}\right],
 \\
 \mathcal R_{Z_\alpha}(\bar\Gamma)_{pqrs}
 &=\sum_t\left[(Z_\alpha)_{pt}\bar\Gamma_{tqrs}
              +(Z_\alpha)_{qt}\bar\Gamma_{ptrs}
 \right.\nonumber\\
 &\hspace{5.5em}\left.
              +(Z_\alpha)_{rt}\bar\Gamma_{pqts}
              +(Z_\alpha)_{st}\bar\Gamma_{pqrt}\right].
 \label{eq:nr-orbital-rdm-response}
\end{align}
The conceptual relaxed full-MO RDMs are
\begin{align}
 \gamma^{\mathrm{rel},\alpha}
 &=\gamma^\alpha+\gamma[\mathbf x_\alpha]
  +\mathcal R_{Z_\alpha}(\bar\gamma),
 \\
 \Gamma^{\mathrm{rel},\alpha}
 &=\Gamma^\alpha+\Gamma[\mathbf x_\alpha]
  +\mathcal R_{Z_\alpha}(\bar\Gamma).
 \label{eq:nr-relaxed-rdms}
\end{align}
The target and CI-transition terms are unweighted.  SA weights enter only
through the barred RDMs in the orbital-response contribution.  The
implementation does not form the full relaxed two-particle RDM in
\cref{eq:nr-relaxed-rdms}; it uses the compact contractions below.

\subsection{Compact hole-space DF weights}

The nuclear-gradient layer differentiates the raw three-center integrals and
metric, avoiding derivatives of the factor $\mathsf X$ in
\cref{eq:nr-orthonormal-df}.  Define the metric-applied three-center tensor
\begin{equation}
 Z^P_{\mu\nu}=\sum_QN_{PQ}J^Q_{\mu\nu}.
 \label{eq:nr-metric-applied-df}
\end{equation}
The auxiliary-indexed tensor $Z^P$ is distinct from the orbital generator
$Z_\alpha$.  Let
\begin{equation}
 C_h=\begin{pmatrix}C_{\core}&C_{\act}\end{pmatrix},
 \qquad
 Z^P_{ij}=\sum_{\mu\nu}(C_h)_{\mu i}
 Z^P_{\mu\nu}(C_h)_{\nu j},
 \quad i,j\in\mathbb H.
 \label{eq:nr-hole-df}
\end{equation}

For any active RDM pair $(\gamma^{\act},\Gamma^{\act})$, define the
spin-free active cumulant
\begin{equation}
 \Lambda^{\act}_{uv,wt}
 =\Gamma^{\act}_{uv,wt}
 -\gamma^{\act}_{uw}\gamma^{\act}_{vt}
 +\frac12\gamma^{\act}_{ut}\gamma^{\act}_{vw}
 \label{eq:nr-active-cumulant}
\end{equation}
and the hole-space one-particle density
\begin{equation}
 \gamma^h=
 \begin{pmatrix}2I_{\core}&0\\0&\gamma^{\act}\end{pmatrix}.
 \label{eq:nr-hole-1rdm}
\end{equation}
The corresponding hole pair density could be written
\begin{equation}
 \Gamma^h_{ij,kl}
 =\gamma^h_{ik}\gamma^h_{jl}
 -\frac12\gamma^h_{il}\gamma^h_{jk}
 +\sum_{uvwt\in\mathbb A}
 \delta_{iu}\delta_{jv}\delta_{kw}\delta_{lt}
 \Lambda^{\act}_{uv,wt},
 \label{eq:nr-hole-2rdm}
\end{equation}
but it is never stored.

The implementation instead builds
\begin{align}
 R^P&=\sum_{ij\in\mathbb H}\gamma^h_{ij}Z^P_{ij},
 \\
 X^P_{ij}
 &=\sum_{kl\in\mathbb H}\gamma^h_{ik}Z^P_{lk}\gamma^h_{lj},
 \\
 L^P_{uw}
 &=\sum_{vt\in\mathbb A}\Lambda^{\act}_{uv,wt}Z^P_{vt}.
 \label{eq:nr-rxl}
\end{align}
The compact three-center and auxiliary-metric derivative weights are
\begin{align}
 \boxed{
 W^{P,h}_{ij}
 =\gamma^h_{ij}R^P-\frac12X^P_{ij}
 +\sum_{uw\in\mathbb A}\delta_{iu}\delta_{jw}L^P_{uw}
 },
 \label{eq:nr-hole-w3}
 \\
 \boxed{
 W_{PQ}
 =-\frac12R^PR^Q
 +\frac14\sum_{ij\in\mathbb H}Z^P_{ij}X^Q_{ij}
 -\frac12\sum_{uw\in\mathbb A}Z^P_{uw}L^Q_{uw}
 }.
 \label{eq:nr-hole-w2}
\end{align}
The $1/2$ and $1/4$ factors are the nonrelativistic spin-free exchange
factors.  These expressions are the factorized evaluation of
\begin{equation}
 W^{P,h}_{ij}=\sum_{kl\in\mathbb H}\Gamma^h_{ik,jl}Z^P_{kl},
 \qquad
 W_{PQ}=-\frac12\sum_{ijkl\in\mathbb H}
 \Gamma^h_{ik,jl}Z^P_{ij}Z^Q_{kl}.
 \label{eq:nr-unfactorized-weights}
\end{equation}
The AO three-center weight is back-transformed only after the hole-space
contractions,
\begin{equation}
 W^P_{\mu\nu}=\sum_{ij\in\mathbb H}
 (C_h)_{\mu i}W^{P,h}_{ij}(C_h)_{\nu j}.
 \label{eq:nr-w3-backtransform}
\end{equation}

\subsection{Target and CI-multiplier contribution}

Define the active base RDMs
\begin{equation}
 \gamma^{\mathrm b}=\gamma^\alpha+\gamma[\mathbf x_\alpha],
 \qquad
 \Gamma^{\mathrm b}=\Gamma^\alpha+\Gamma[\mathbf x_\alpha].
 \label{eq:nr-base-rdms}
\end{equation}
Because the projected CI multiplier satisfies $s[\mathbf x_\alpha]=0$, the
closed-core transition contribution vanishes.  The base hole 1-RDM therefore
retains the $2I_{\core}$ block in \cref{eq:nr-hole-1rdm}, with only its active
block replaced by $\gamma^{\mathrm b}$.  Applying
\cref{eq:nr-active-cumulant,eq:nr-rxl,eq:nr-hole-w3,eq:nr-hole-w2} to the base
RDMs defines
\begin{equation}
 R^{P,\mathrm b},\quad X^{P,\mathrm b},\quad L^{P,\mathrm b},
 \quad W^{P,h,\mathrm b},\quad W^{\mathrm b}_{PQ}.
\end{equation}

\subsection{Orbital-multiplier contribution}

Apply the same compact equations to $(\bar\gamma,\bar\Gamma)$ to obtain
$\bar R$, $\bar X$, $\bar L$, $\bar W^{P,h}$, and $\bar W_{PQ}$.  With
$\dot C=CZ_\alpha$ and $\dot C_h$ its core and active columns, the transformed
DF response at fixed AO $Z^P_{\mu\nu}$ is
\begin{equation}
 \dot Z^P_{ij}
 =\sum_{\mu\nu}\left[
 (\dot C_h)_{\mu i}Z^P_{\mu\nu}(C_h)_{\nu j}
 +(C_h)_{\mu i}Z^P_{\mu\nu}(\dot C_h)_{\nu j}
 \right].
 \label{eq:nr-hole-df-response}
\end{equation}
Holding the barred RDMs fixed gives
\begin{align}
 \dot{\bar R}^{P}
 &=\sum_{ij}\bar\gamma^h_{ij}\dot Z^P_{ij},
 \\
 \dot{\bar X}^{P}_{ij}
 &=\sum_{kl}\bar\gamma^h_{ik}\dot Z^P_{lk}\bar\gamma^h_{lj},
 \\
 \dot{\bar L}^{P}_{uw}
 &=\sum_{vt}\bar\Lambda^{\act}_{uv,wt}\dot Z^P_{vt},
 \\
 \dot{\bar W}^{P,h}_{ij}
 &=\bar\gamma^h_{ij}\dot{\bar R}^{P}
 -\frac12\dot{\bar X}^{P}_{ij}
 +\sum_{uw\in\mathbb A}\delta_{iu}\delta_{jw}
  \dot{\bar L}^{P}_{uw}.
 \label{eq:nr-directional-w3}
\end{align}
Differentiating the metric weight gives
\begin{align}
 \dot{\bar W}_{PQ}
 &=-\frac12\left(\dot{\bar R}^{P}\bar R^Q
                  +\bar R^P\dot{\bar R}^{Q}\right)
 \nonumber\\
 &\quad+\frac14\sum_{ij}\left(
 \dot Z^P_{ij}\bar X^Q_{ij}+Z^P_{ij}\dot{\bar X}^Q_{ij}\right)
 \nonumber\\
 &\quad-\frac12\sum_{uw}\left(
 \dot Z^P_{uw}\bar L^Q_{uw}+Z^P_{uw}\dot{\bar L}^Q_{uw}\right).
 \label{eq:nr-directional-w2}
\end{align}

\subsection{Arrays passed to the AO derivative driver}

The relaxed AO one-particle density is
\begin{align}
 D^{\mathrm{rel},\alpha}_{\mu\nu}
 &=2\sum_{m\in\mathbb C}C_{\mu m}C_{\nu m}
 +\sum_{uv\in\mathbb A}C_{\mu u}\gamma^{\mathrm b}_{uv}C_{\nu v}
 \nonumber\\
 &\quad+2\sum_{m\in\mathbb C}\left(
 \dot C_{\mu m}C_{\nu m}+C_{\mu m}\dot C_{\nu m}\right)
 \nonumber\\
 &\quad+\sum_{uv\in\mathbb A}\left(
 \dot C_{\mu u}\bar\gamma_{uv}C_{\nu v}
 +C_{\mu u}\bar\gamma_{uv}\dot C_{\nu v}\right).
 \label{eq:nr-relaxed-ao-1rdm}
\end{align}
The relaxed three-center weight, called \texttt{W3} in the implementation, is
\begin{align}
 W^{P,\alpha}_{\mu\nu}
 &=\sum_{ij}(C_h)_{\mu i}W^{P,h,\mathrm b}_{ij}(C_h)_{\nu j}
 \nonumber\\
 &\quad+\sum_{ij}\left[
 (\dot C_h)_{\mu i}\bar W^{P,h}_{ij}(C_h)_{\nu j}
 +(C_h)_{\mu i}\bar W^{P,h}_{ij}(\dot C_h)_{\nu j}
 \right]
 \nonumber\\
 &\quad+\sum_{ij}(C_h)_{\mu i}
 \dot{\bar W}^{P,h}_{ij}(C_h)_{\nu j}.
 \label{eq:nr-relaxed-w3}
\end{align}
The relaxed metric weight, called \texttt{W2}, is
\begin{equation}
 W^\alpha_{PQ}=W^{\mathrm b}_{PQ}+\dot{\bar W}_{PQ}.
 \label{eq:nr-relaxed-w2}
\end{equation}
The response of both outer $C_h$ factors and of the internally transformed
$Z^P_{ij}$ is therefore included exactly once.

The common AO derivative-integral layer evaluates the two-electron part as
\begin{equation}
 E^x_{2,\alpha}
 =\sum_{P\mu\nu}W^{P,\alpha}_{\mu\nu}(J^P_{\mu\nu})^x
 +\sum_{PQ}W^\alpha_{PQ}(M_{PQ})^x.
 \label{eq:nr-df-gradient-driver}
\end{equation}
Combining this with the one-electron, overlap, and nuclear-repulsion terms
gives the implemented gradient:
\begin{equation}
 \boxed{
 E_\alpha^x
 =V_{\mathrm{NN}}^x
 +\sum_{\mu\nu}D^{\mathrm{rel},\alpha}_{\mu\nu}h^x_{\mu\nu}
 -\sum_{\mu\nu}W^{S,\alpha}_{\mu\nu}S^x_{\mu\nu}
 +\sum_{P\mu\nu}W^{P,\alpha}_{\mu\nu}(J^P_{\mu\nu})^x
 +\sum_{PQ}W^\alpha_{PQ}(M_{PQ})^x
 }.
 \label{eq:nr-final-gradient}
\end{equation}

\section{Implementation summary and scaling}

The response solver is matrix free and applies all four orbital/CI blocks
without storing any dense response matrix.

The minimum-storage DF choices used by the implementation are:
\begin{itemize}
 \item keep only $B^P_{pu}$ rather than a general--general transformed
       three-index tensor;
 \item evaluate transition exchange through signed low-rank density factors;
 \item contract active two-particle RDMs directly with general--active DF
       blocks;
 \item form only the $O(N_{\mathrm A}^4)$ active integral response needed by
       the existing CI sigma builder;
 \item apply the coupled response matrix with GMRES rather than storing the
       orbital/CI Hessian; and
 \item assemble nuclear derivative weights through the hole-space $R$, $X$,
       and $L$ contractions without storing an $O(N_{\mathrm H}^4)$ pair
       density.
\end{itemize}
The persistent transformed response workspace is
$O(N_{\mathrm{aux}}N_{\mathrm{MO}}N_{\mathrm A})$.  Gradient assembly needs
$O(N_{\mathrm{aux}}N_{\mathrm H}^2)$ compact work storage in addition to the
AO three-center and auxiliary-metric derivative weights required by the
common DF gradient driver.

\end{document}
